\section{Implemented Synthetic Time-Series Generator}
\label{sec:implemented-generator}
%
% PURPOSE OF THIS CHAPTER
% Explain the generator that was actually implemented, separating implemented
% behavior from the broader design proposed in Chapter 3.  Connect every design
% choice to a source file, public object, test, or executed notebook.
%
% \subsection{Scope and implemented pipeline}
% - Begin with one compact pipeline diagram:
%   user formula or sampled tree -> validated AST -> executable recurrence ->
%   one or more trajectories -> canonical batch.
% - State the implemented operator vocabulary and deliberately deferred features.
% - Distinguish a mechanism (one retained AST) from a trajectory (one realization
%   generated from that mechanism).
%
% \subsection{Temporal convention}
% - Define the forecast origin, the meaning of t, and the convention that
%   LAG(j, ell) represents x_{j,t-ell} with ell >= 1.
% - Explain how context and future indices are represented in arrays and plots.
% - Give one hand-computed recurrence example to remove any indexing ambiguity.
%
% \subsection{Abstract syntax tree and fixed grammar}
% - Define AST at first use and show the implemented node types.
% - Explain terminals such as constants, time t, and LAG(j, ell), and operators
%   such as addition and multiplication.
% - Reuse a small tree diagram and show its equivalent infix formula and
%   recurrence.
% - Explain that the grammar is fixed by the benchmark author in the first
%   version rather than supplied by each library user.
%
% \subsection{Formula input, parsing, and representations}
% - Show the syntax accepted from a user and how instructions are exposed by the
%   library object.
% - Describe conversion among text, AST, executable function, and rendered
%   mathematical formula.
% - Discuss canonical serialization and how commutative children are handled;
%   do not claim full symbolic equivalence unless it has been implemented and
%   tested.
%
% \subsection{Constrained random unary-binary tree sampling}
% - Present the sampler pseudocode and cite the exact implementation file.
% - Explain zero-operator trees, unary nodes, binary nodes, terminals, maximum
%   depth, lags, series count, coefficients, and operator probabilities.
% - Explain how one root seed is deterministically split into mechanism and
%   trajectory seeds.
% - Compare the implemented sampler with the expression-first precedents from
%   Section 1.2, especially D'Ascoli et al. and Meznar et al.; cite only claims
%   supported by the corresponding reviewed sources.
%
% \subsection{Executing a mechanism and generating trajectories}
% - Explain initial values, recurrence execution, optional innovation terms,
%   context length, forecast horizon, and the relationship between total length
%   and returned windows.
% - Show why several trajectories can share one formula while differing through
%   their trajectory-level random draws.
% - Describe current finite-value or pilot checks accurately.  Treat formal
%   stability guarantees as future work unless they are implemented.
%
% \subsection{Jobs, parallel generation, and canonical batches}
% - Define one job as one sampled mechanism plus its requested trajectories.
% - Explain the minimum number of distinct mechanisms required in a batch.
% - Describe the PyTorch-style batch-generation objective without claiming an
%   IterableDataset/DataLoader integration before it exists.
% - Explain that output ordering and seed derivation are independent of worker
%   scheduling, so changing n_jobs must preserve the batch content.
% - Report the relevant hash/reproducibility tests in the test suite rather than
%   adding another validation experiment to the notebook.
%
% \subsection{User-facing visualization and executed examples}
% - Describe the frontend interaction: sample a tree, inspect its equation, and
%   draw several trajectories from the retained mechanism.
% - Summarize notebooks 1--3 without reproducing all notebook prose.
% - Include the mechanism-by-column and trajectory-by-row visualization when it
%   becomes a stable report figure.
%
% \subsection{Implementation validation and current limitations}
% - Summarize parser, serialization, reproducibility, parallelism, and batch
%   tests with exact counts only when the final experiment snapshot is frozen.
% - Identify unsupported operators, multivariate behavior, stochastic families,
%   stability certification, and performance constraints from repository
%   evidence at writing time.
% - End by stating which retained generator information is consumed by the
%   explanation-ground-truth layer in Chapter 5.
